% --- Back cover ---
% Large PERSONIX wordmark, no logo image. Cream paper inherited globally.
% Absolute (tikz overlay) positioning so the footer sticks to the bottom.
\clearpage
\thispagestyle{empty}
\begin{tikzpicture}[remember picture, overlay]
  % --- Wordmark ---
  \node[anchor=north, font=\sffamily\bfseries\fontsize{48}{52}\selectfont]
    at ([yshift=-26mm]current page.north) {PERSONIX};
  \node[anchor=north, font=\rmfamily\itshape\large]
    at ([yshift=-44mm]current page.north) {Ασυμβίβαστη Αλλαγή};
  \node[anchor=north, font=\rmfamily\small,
        text width=\paperwidth-50mm, align=center]
    at ([yshift=-54mm]current page.north)
    {Ένα Μη Λογοκρίσιμο και Αδιάφθορο Αποκεντρωμένο Δίκτυο Φήμης};
  % --- Body ---
  \node[anchor=north, text width=\paperwidth-50mm, align=justify, font=\small]
    at ([yshift=-72mm]current page.north)
    {%
      Το κράτος λειτουργούσε όταν η επικοινωνία ήταν αργή, οι αποφάσεις ήταν
      τοπικές, και η εξουσία ζούσε στην ίδια πόλη που κυβερνούσε. Τα σημερινά
      δίκτυα αντιστρέφουν και τις τρεις παραδοχές --- και οι θεσμοί που χτίστηκαν
      πάνω τους δεν ταιριάζουν πια στον κόσμο που κυβερνούν. Αυτό το βιβλίο
      περιγράφει ένα εργαλείο που ταιριάζει: ένα αποκεντρωμένο δίκτυο φήμης όπου
      η ταυτότητα είναι δική σου να την κρατήσεις, η αλήθεια είναι δημόσια
      ελέγξιμη, και η δωροδοκία είναι αυτοκτονία για τη φήμη.

      \smallskip
      Ούτε μανιφέστο. Ούτε πρόβλεψη. Ένα λειτουργικό σχέδιο για το πώς μπορεί
      να χτιστεί ο διάδοχος του κράτους --- εθελοντικά, μια δήλωση τη φορά.
    };
  % --- Footer (anchored to the bottom of the page) ---
  \node[anchor=south, font=\sffamily\footnotesize,
        text width=\paperwidth-50mm, align=center]
    at ([yshift=12mm]current page.south)
    {%
      Pavel Kudrna\quad\textbullet\quad Πράγα, 2026\par
      \href{https://personix.org}{personix.org}\par
      Αδειοδοτημένο υπό την άδεια Creative Commons Attribution-ShareAlike 4.0
      International (CC BY-SA 4.0).
    };
\end{tikzpicture}%
\mbox{}%
\clearpage
